% SPDX-License-Identifier: Apache-2.0
% Copyright (c) 2026 finevals.ai and contributors.
%
% Credit Foundation Models -- paper draft.
% Neutral single-column article preamble; the body ports unchanged if a specific
% template is adopted later.
%
% Build:  tectonic paper.tex      (or: pdflatex paper.tex, run twice for references)
\documentclass[10pt,letterpaper]{article}
\usepackage[margin=1in]{geometry}
\usepackage{booktabs}
\usepackage{amsmath,amssymb}
\usepackage{graphicx}
\usepackage{url}
\usepackage[hidelinks]{hyperref}
\usepackage{caption}
\captionsetup{font=small}

\title{\textbf{Feed the Model Before You Grow It:\\
A Sequence Foundation Model for Credit Risk,\\
Evaluated Out-of-Time}}

\author{%
  Sriram Krishnan \\
  \normalsize finevals.ai \\
  \normalsize \texttt{sriram.arun@gmail.com}
}
\date{}

\begin{document}
\maketitle

\begin{abstract}
Credit risk models in production are overwhelmingly \emph{point-in-time}: gradient-boosted trees
applied to a single feature row summarising a loan as of today. We ask whether a \emph{sequence}
foundation model, pretrained without labels on the full month-by-month repayment history of
millions of loans, produces a better risk signal --- and we insist on measuring it the way
deployment actually tests a model: trained on the past, evaluated on future years the model has
never seen. On a public panel of US single-family mortgage performance data (2000--2024,
$\approx$3.3B loan-month observations), an encoder-only masked-language-model with a three-branch
architecture over key-value-time tokenized credit events reaches \textbf{0.8468 ROC-AUC / 0.0175
average precision} on a calendar out-of-time window, against \textbf{0.7913 / 0.0057} for a
strong leakage-audited gradient-boosting baseline on the identical population --- \textbf{+0.056
ROC and $3.1\times$ AP} at a $0.14\%$ default base rate. Beyond the headline we report a
controlled scaling decomposition that we believe is the more useful contribution. Growing the
backbone $2.6\times$ on unchanged data changes nothing; growing it a further $4\times$ on a
matched corpus changes nothing either --- a $25.7$M model reaches $0.8488$ ROC against the $100.9$M
model's $0.8468$. Data accounts for essentially the whole gain. Two null results at two scales give
one practical rule: \emph{feed the model before you grow it}. We further report an honest negative result on
prepayment, a crisis-regime stress test, and an end-to-end independent reproduction of the
headline on a refactored codebase. The complete training framework is released under Apache~2.0.
\end{abstract}

\section{Introduction}
\label{sec:intro}

A bank deciding whether a borrower will default next year has, in principle, a rich object to
reason over: the loan's entire repayment history --- every monthly balance, every delinquency
transition, every modification --- stretching back to origination. In practice, almost all
deployed credit models collapse that history into a single row of engineered features and hand it
to a gradient-boosted tree. The sequence is discarded before modelling begins.

This paper asks a narrow, empirical question: \emph{does keeping the sequence pay, and if so, what
is doing the work?} We treat a loan's history as a sentence in a domain-specific language and
pretrain a transformer encoder on it with a masked-token objective, no labels required. The
pretrained backbone is then adapted to specific tasks (default, prepayment) with a lightweight
head.

Two commitments shape everything that follows.

\textbf{Evaluation must be out-of-time.} Credit models fail when the world moves, not when a random
row is held out. We therefore fit on observation snapshots from one set of calendar years and test
on later ones whose outcome windows fall entirely after the training data ends. Random splits, and
even loan-disjoint random splits, systematically flatter a credit model; we avoid both.

\textbf{Scale claims must be decomposed.} It is easy to report that a bigger model scored better and
imply that capacity was responsible. We ran the controls: a capacity-only arm, a data-only arm, and
the combined arm. The capacity-only arm is a null result, and we report it prominently --- it is
what makes the positive results credible.

Our contributions are:
\begin{enumerate}
  \item A sequence foundation model for tabular credit panels that beats a leakage-audited
        gradient-boosting baseline out-of-time by $+0.056$ ROC-AUC and $3.1\times$ average
        precision at a $0.14\%$ base rate (Section~\ref{sec:results}).
  \item A controlled scaling decomposition separating backbone capacity from corpus size, yielding
        a concrete practitioner rule (Section~\ref{sec:scaling}).
  \item Honest negative and stress results: prepayment is near-unpredictable from borrower history
        alone; the benign in-period window favours features; a crisis-regime transfer test
        (Section~\ref{sec:stress}).
  \item An end-to-end independent reproduction of the headline on a refactored codebase, and the
        complete training framework released under Apache~2.0 (Section~\ref{sec:repro}).
\end{enumerate}

\section{Related Work}
\label{sec:related}

\textbf{Foundation models for tabular and event data.}
Transformer architectures adapted to tabular inputs --- attention over feature embeddings rather
than tokens --- have been explored extensively, and a parallel line treats sequences of discrete
events (transactions, clicks, clinical codes) as a language amenable to masked pretraining. Our
setting sits at the intersection: the unit of study is a \emph{panel}, one row per entity per
month, mixing static origination attributes with a long dynamic sequence.

\textbf{Encoder-only masked pretraining.}
We adopt the bidirectional masked-token objective of BERT~\cite{devlin2019bert} rather than a
causal objective. The task is to \emph{score} a completed history, not to generate one; bidirectional
context is available at inference and there is no generation surface to hallucinate.

\textbf{Financial event-sequence foundation models.}
The closest published work is PRAGMA~\cite{pragma2026}, a family of foundation models for
multi-source \emph{banking event} sequences (10M--1B parameters). Its design choices are close to
ours and were arrived at independently for the same reasons: an encoder-only bidirectional backbone
chosen because the goal is transfer rather than generation, separate encoder branches for profile
state and events whose outputs are fused by a history encoder, and adaptation by embedding probe or
LoRA. PRAGMA reports large relative gains on credit scoring ($+130.2\%$ PR-AUC, $+12.4\%$ ROC-AUC
over a task-specific baseline).

Two differences matter for what a reader can check. First, the substrate: PRAGMA models
heterogeneous banking events (transactions, product interactions) from a proprietary corpus,
whereas we model a \emph{loan-performance panel} --- one row per loan per month, static origination
attributes plus a long monthly sequence --- from a public source. Second, and more consequentially,
PRAGMA states that it reports \emph{only relative improvements, as absolute metrics are
commercially sensitive}. We report absolute ROC-AUC and average precision against a disclosed
baseline on public data, with the code released, so every number here can be re-derived rather than
taken on trust. An industrial transaction foundation model blueprint~\cite{nvidiatfm} provides our
engineering baseline and the convention of judging average precision first.

\textbf{Scaling laws.}
Compute-optimal scaling~\cite{hoffmann2022chinchilla} predicts that model size and data must grow
together. Our decomposition (Section~\ref{sec:scaling}) is a direct, small-scale test of that
prediction in a domain where the data ceiling is a real operational constraint, and it reproduces
the qualitative prediction: capacity added without data is wasted.

\textbf{Parameter-efficient adaptation.}
We compare frozen-embedding, low-rank adaptation~\cite{hu2021lora}, and full fine-tuning as
interchangeable heads on one backbone.

\section{Method}
\label{sec:method}

\subsection{From panel rows to a language}
The raw corpus is a panel: one row per loan per month, combining \emph{static} attributes fixed at
origination (term, loan-to-value, debt-to-income, credit score band, occupancy, property type) with
\emph{dynamic} monthly state (current balance, loan age, remaining term, delinquency status).

We serialise each loan into a token sequence with a key-value-time (KVT) scheme:
\begin{itemize}
  \item Every field becomes a single fused token \texttt{field=value}, so the model never has to
        learn which column a bare value came from.
  \item Continuous fields are quantile-binned, with bin boundaries \emph{forced} at points where
        credit behaviour genuinely changes (loan-to-value at 80/90/95/97, debt-to-income at 43).
        These are regulatory and pricing cliffs; letting a quantiser smear across them destroys
        signal that underwriters treat as categorical.
  \item Each monthly block is anchored by a relative time token \texttt{t=}$\langle$age$\rangle$ and
        an absolute calendar token \texttt{cal=}$\langle$YYYYQ$n\rangle$. The calendar token is what
        allows the model to represent macro regime --- a 2009 month and a 2021 month are not
        exchangeable.
\end{itemize}
The vocabulary is small (552 tokens) and is fit \emph{on the training split only}, then frozen.

\subsection{Three-branch encoder}
Static attributes and monthly events have different statistics and different roles, so we encode
them separately before fusing:
\[
\underbrace{\text{static fields}}_{\text{Profile encoder (3L)}} \;\;\oplus\;\;
\underbrace{\text{monthly events}}_{\text{Event encoder (5L)}}
\;\longrightarrow\;
\underbrace{\text{History encoder (6L)}}_{\text{sequence fusion}}
\;\longrightarrow\; \texttt{[USR]} \text{ loan embedding}
\]
The Profile encoder sees the origination record; the Event encoder contextualises each month's
fields within that month; the History encoder attends across months. A distinguished
\texttt{[USR]} position aggregates the loan into a single vector used by downstream heads.
Blocks use RoPE positional encoding, RMSNorm, SwiGLU feed-forwards, and scaled dot-product
(Flash) attention. We report two sizes: $25.7$M parameters at $d{=}384$ and $100.9$M at $d{=}768$,
identical in depth and differing only in width.

\subsection{Pretraining objective}
We minimise a masked-token cross-entropy with three complementary corruption sources, applied
jointly: (i) $15\%$ of individual tokens, (ii) $10\%$ of whole monthly event blocks, and (iii)
$10\%$ of entire field \emph{types} across the sequence. Masking whole events forces temporal
inference (reconstruct a month from its neighbours) and masking a field type forces cross-field
inference (reconstruct balance from delinquency and age). Selected positions follow the standard
80/10/10 replace/random/keep policy. Formally, for masked index set $\mathcal{M}$,
\[
\mathcal{L} = -\frac{1}{|\mathcal{M}|}\sum_{i\in\mathcal{M}} \log p_\theta\!\left(x_i \mid x_{\setminus \mathcal{M}}\right).
\]

\subsection{Adaptation}
The pretrained backbone is adapted with one of three interchangeable heads: \emph{frozen}
(embeddings fed to a classifier, backbone untouched), \emph{LoRA}~\cite{hu2021lora} (low-rank
updates, $r{=}8$), or \emph{full} fine-tuning. Tasks are declared in configuration rather than
code: a task specifies an event column, a horizon, and a gate (e.g.\ observe only loans performing
at the cutoff), so adding a task requires no new modelling code.

\section{Experimental Setup}
\label{sec:setup}

\subsection{Corpus}
We use a public panel of US single-family mortgage performance data covering 2000--2024:
approximately $3.3$B loan-month observations over tens of millions of fixed-rate loans, with 113
published fields.\footnote{We describe the corpus by its properties rather than by name. It is a
widely used public loan-performance release; the released code contains the ingestion adapter and
schema needed to reproduce every result once the reader supplies the source files.} Pretraining
uses deterministic loan-hash samples ($4\%$ and $10\%$) of the book; a $4\%$ sample was verified
representative against the full population (pooled default rate $0.671\%$ vs $0.648\%$).

\subsection{Leakage discipline}
Credit datasets punish carelessness, so the protocol is enforced mechanically rather than by
convention:
\begin{itemize}
  \item Splits are by \emph{loan}, never by row, and ordered temporally by origination.
  \item Outcome and contemporaneous-state columns (current delinquency, zero-balance code,
        foreclosure and disposition dates, loss amounts) are excluded from features for
        \emph{both} the
        foundation model and the baseline; the baseline additionally gates to loans performing at
        the observation date.
  \item Tokenizer vocabulary and bin edges are fit on the training split only.
  \item Evaluation is calendar out-of-time with loan-disjoint and embargo guards.
\end{itemize}
Each pipeline stage ships an artifact validator that re-derives its own output and is tested to
\emph{fail} on deliberately corrupted input.

\subsection{Protocol and metrics}
The main protocol fits a head on observation snapshots from 2016--2021 and tests on 2022 and 2023
snapshots, whose 12-month default windows fall in 2023--2024 --- years neither the backbone nor the
head has seen. The test population is $1{,}782{,}453$ loan observations at a $0.14\%$ default rate.

We report ROC-AUC and average precision (AP, equivalently PR-AUC). At a $0.14\%$ base rate we
regard \textbf{AP as the operational metric}: it measures precision in the risky tail, which is
where a credit decision is actually made. ROC is reported for comparability but is generous at
extreme class imbalance.

\subsection{Baseline}
The baseline is a gradient-boosted tree over 57 engineered, leakage-audited features on the
identical observation population and identical label definition. It is intended to be strong and
honest rather than a straw man: it uses the same gating and the same excluded-column list.

\section{Results}
\label{sec:results}

\subsection{Main out-of-time result}
\begin{table}[h]
\centering
\begin{tabular}{lcc}
\toprule
\textbf{Model} & \textbf{ROC-AUC} & \textbf{AP} \\
\midrule
Gradient-boosting baseline (57 features)      & 0.7913 & 0.0057 \\
Credit FM, 25.7M parameters (4\% corpus)      & 0.8257 & 0.0113 \\
\textbf{Credit FM, 25.7M (10\% corpus)}       & \textbf{0.8488} & 0.0156 \\
Credit FM, 100.9M (10\% corpus)               & 0.8468 & \textbf{0.0175} \\
\bottomrule
\end{tabular}
\caption{Calendar out-of-time performance. Head fit on 2016--2021 snapshots; tested on 2022--2023
snapshots with outcomes in 2023--2024. Test population $1{,}782{,}453$ loans, $0.14\%$ default.}
\label{tab:main}
\end{table}

The foundation model wins on both metrics, and the margin that matters operationally is the
$3.1\times$ improvement in average precision.

\subsection{Adaptation mode matters}
At the 25.7M scale, holding data and protocol fixed: frozen embeddings reach $0.7309$, LoRA
$0.8068$, and full fine-tuning $0.8257$. Frozen embeddings alone do \emph{not} beat good features
--- consistent with prior reports --- and the pretrained representation only pays once the backbone
is allowed to adapt.

\subsection{Scaling: what actually moves the needle}
\label{sec:scaling}

\begin{table}[h]
\centering
\begin{tabular}{lllcc}
\toprule
\textbf{Arm} & \textbf{Backbone} & \textbf{Pretrain / fine-tune data} & \textbf{ROC-AUC} & \textbf{AP} \\
\midrule
Reference               & 25.7M  & 4\% / 4\%    & 0.8257 & 0.0113 \\
Capacity only           & 66.8M  & 4\% / 4\%    & 0.8223 & $\approx$ ref \\
Fine-tune data only     & 25.7M  & 4\% / 10\%   & 0.8406 & 0.0145 \\
\textbf{Pretrain data too} & \textbf{25.7M} & \textbf{10\% / 10\%} & \textbf{0.8488} & 0.0156 \\
And capacity            & 100.9M & 10\% / 10\%  & 0.8468 & \textbf{0.0175} \\
\bottomrule
\end{tabular}
\caption{Scaling decomposition, one variable at a time. The last two rows share an identical corpus
and an identical recipe (effective batch 256, 20k steps) and differ only in backbone width, which
isolates capacity. Test population is identical across the final three rows.}
\label{tab:scaling}
\end{table}

Four findings, one variable at a time:
\begin{enumerate}
  \item \textbf{Capacity alone does nothing.} A $2.6\times$ larger backbone on unchanged data is
        flat ($0.8223$ vs $0.8257$, within noise). The smaller model was already data-bound.
  \item \textbf{Data alone recovers most of the gain.} Holding the backbone fixed and growing the
        data captures $+0.0149$ of the $+0.0211$ total ROC gain ($\approx 71\%$), and $+0.0032$ AP.
  \item \textbf{Capacity still does nothing, even once the data is there.} Holding the corpus and
        the recipe fixed and widening $25.7$M~$\rightarrow$~$100.9$M moves ROC by $-0.0020$ and AP by
        $+0.0019$ --- both inside one standard error. The 25.7M model reaches $0.8488$ ROC against
        the 100.9M model's $0.8468$, at a quarter of the parameters.
\end{enumerate}

On ROC, essentially all of the gain is data: the matched-data 25.7M arm alone is $+0.0231$ over the
reference, exceeding the $+0.0211$ of the largest model. On AP, data contributes $\approx$69\% of the
total and capacity $\approx$31\%. At this data scale the two backbones are statistically
indistinguishable, and the efficient one is four times cheaper.

The practitioner's rule: \textbf{feed the model before you grow it}. We regard the two null results
--- capacity added at $4\%$ data, and capacity added at $10\%$ data --- as the most transferable
finding in the paper.

\subsection{Stress and negative results}
\label{sec:stress}
\textbf{Crisis regime.} Trained on originations through 2006 and tested on 2008--2010 defaults ---
a genuine distribution shift --- the model reaches $0.7819$ ROC / $0.0248$ AP against
$0.757$ / $0.024$ for the baseline. The margin narrows under severe shift but does not invert.

\textbf{Benign in-period window.} On a same-period window with no regime shift, engineered features
win narrowly. This is the expected and honest result: the sequence model earns its keep precisely
when the world moves.

\textbf{Prepayment (negative).} The identical pipeline, retargeted to 12-month prepayment by
changing one configuration line, reaches only $0.6259$ ROC. Prepayment is dominated by exogenous
interest-rate dynamics that a borrower-history model cannot see. We report this because a framework
that only ever produces good numbers should not be trusted.

\subsection{Independent reproduction}
\label{sec:repro}
After nine refactoring pull requests to the codebase, we re-ran the headline end to end from raw
data on fresh output paths. Deterministic stages matched exactly (identical split loan counts,
$254.4$M training rows, $2{,}997{,}726{,}396$ tokens, and an identical $1{,}782{,}453$-observation
evaluation population). The stochastic endpoint landed at $0.8447$ ROC / $0.0160$ AP against the
original $0.8468$ / $0.0175$ --- inside the pre-registered run-to-run band. The reproduction also
surfaced one genuine defect (a schema-unification bug across per-quarter shards), now fixed and
regression-tested.

\section{Discussion}
\label{sec:discussion}

\textbf{Why sequences win out-of-time and not in-period.} A snapshot feature vector and a full
history contain similar information about a stable world. They diverge when the world shifts: the
trajectory (how quickly a borrower cured a delinquency, how balance moved through a rate cycle)
generalises across regimes in a way that a level does not. Our benign-window result and our
out-of-time result together isolate exactly this.

\textbf{Why AP and not ROC.} At $0.14\%$ positives, ROC differences are compressed and easy to
over-read. AP measures the quantity a credit officer experiences --- of the loans flagged riskiest,
how many truly default --- and it is where our margin is largest ($3.1\times$).

\textbf{Implications for practitioners.} The scaling decomposition suggests the first move for an
institution with a modest labelled book is not a larger model but a larger \emph{unlabelled} pretraining
corpus, which most lenders already possess and do not exploit.

\section{Limitations}
\label{sec:limitations}
\begin{itemize}
  \item \textbf{Single corpus.} All headline results come from one national mortgage market. A
        second asset class (business-to-business invoice payment behaviour) has been onboarded
        through the same interface but is not yet a published result.
  \item \textbf{Benign test years.} 2023--2024 were low-default years; absolute AP is small for all
        models and the \emph{relative} lift is the claim, not the absolute value.
  \item \textbf{Statistical power.} With $\approx$2{,}500 test positives, a ROC increment of
        $\pm0.01$ is roughly one standard error. The headline margin over the baseline ($+0.056$)
        is far outside that band; the smaller within-family increments are correspondingly weaker
        evidence and are labelled as such.
  \item \textbf{Test-population note.} Arms trained on different sample fractions are evaluated on
        the corresponding deterministic loan-hash test populations; the $4\%$ sample is a subset of
        the $10\%$ by construction, so the populations match in distribution but are not row-identical
        across all pairs.
  \item \textbf{Scale ceiling.} We test to $100.9$M parameters and $3.0$B tokens. The full corpus
        supports roughly an order of magnitude more; that headroom is untested.
  \item \textbf{Corpus naming.} We characterise the dataset by its properties rather than naming the
        publisher. The released ingestion adapter and schema make the mapping unambiguous for a
        reader who obtains the public source files.
\end{itemize}

\section{Conclusion}
Treating a loan's repayment history as a sequence, and pretraining on it without labels, produces a
materially better credit risk signal than a strong point-in-time baseline --- and the advantage
appears exactly where it matters, on future years the model has never seen. The scaling
decomposition matters more than the headline: capacity added to a data-bound model bought nothing,
while data bought most of the gain. The framework, the reference implementation, and the evaluation
protocol are released under Apache~2.0 so that both the positive and the negative results can be
checked.

\section*{Reproducibility Statement}
The complete training framework, configuration recipes, per-stage artifact validators, and the
runnable pipeline are released under Apache~2.0. Every headline number in this paper corresponds to
a committed configuration and a generated report file. Section~\ref{sec:repro} documents an
independent end-to-end re-run of the headline result on a refactored codebase.

\sloppy
\begin{thebibliography}{9}
\bibitem{devlin2019bert}
J.~Devlin, M.-W. Chang, K.~Lee, and K.~Toutanova.
\newblock BERT: Pre-training of Deep Bidirectional Transformers for Language Understanding.
\newblock In \emph{NAACL-HLT}, 2019.

\bibitem{hu2021lora}
E.~Hu, Y.~Shen, P.~Wallis, Z.~Allen-Zhu, Y.~Li, S.~Wang, L.~Wang, and W.~Chen.
\newblock LoRA: Low-Rank Adaptation of Large Language Models.
\newblock \emph{arXiv:2106.09685}, 2021.

\bibitem{hoffmann2022chinchilla}
J.~Hoffmann et al.
\newblock Training Compute-Optimal Large Language Models.
\newblock \emph{arXiv:2203.15556}, 2022.

\bibitem{pragma2026}
M.~Ostroukhov, R.~Mikhailov, V.~Iashin, A.~Sokolov, A.~Akshonov, V.~Protasov, D.~Beloborodov,
V.~Mullin, R.~Y. Enzmann, G.~Kolovos, J.~Renders, P.~Nesterov, and A.~Repushko.
\newblock PRAGMA: Revolut Foundation Model.
\newblock \emph{arXiv:2604.08649}, 2026. Revolut Research and NVIDIA.

\bibitem{nvidiatfm}
NVIDIA AI Blueprints.
\newblock Transaction Foundation Model blueprint.
\newblock \url{https://github.com/NVIDIA-AI-Blueprints/transaction-foundation-model}.
\end{thebibliography}

\end{document}
